\chapter[Cores and galaxies]{Regular Cores and Galactic Backreaction}
\label{ch:cores-galaxies}

\paragraph{Learning outcomes.}
After this chapter the reader should be able to:
\begin{enumerate}
  \item write the static, spherically symmetric reduction of the
        infrared equation
        \(G_{\mu\nu}+\Lambda_{\mathrm{eff}}g_{\mu\nu}
        =8\pi\Geff T_{\mu\nu}^{\mathrm{total}}\)
        with the Einsteinian cubic term included;
  \item state the regular-core profile
        \(f(r)=1-2Mr^2/(r^3+r_0^3)-\Lambda_{\mathrm{eff}}r^2/3\)
        as a phenomenological model, and state the scope of the
        no-go on its exactness (Section~\ref{sec:status-J});
  \item reduce the reversible part of \(T_{\mu\nu}^{\mathrm{ent}}\)
        to a modification of Poisson's equation;
  \item derive a schematic circular-velocity formula and list the
        observations that would distinguish it from a particulate
        halo;
  \item enter Chapter~14 knowing that any time-dependent dissipative
        term would be a separate addition to the static extra
        potential used here.
\end{enumerate}

The geometric programme of Parts~II--IV is applied to two symmetric
sectors: static spherical interiors and weak-field galactic
exteriors. The regular-core profile below remains a model conditional
on Problem~J, not a consequence of the earlier theorems.

\section{Static spherical ansatz}
\label{sec:ch13-ansatz}

Take the emergent metric in Schwarzschild coordinates,
\begin{equation}
  \label{eq:sss}
  ds^2
  =-f(r)\,dt^2
   +\frac{dr^2}{f(r)}
   +r^2 d\Omega^2,
\end{equation}
and restrict \(T_{\mu\nu}^{\mathrm{total}}\) to a static, spherical
form. The Einstein tensor of~\eqref{eq:sss} has the standard
components
\begin{equation}
  \label{eq:G-sss}
  G^t{}_t=G^r{}_r
  =-\frac{1-f-rf'}{r^2},
  \qquad
  G^\theta{}_\theta
  =-\frac{rf''+2f'}{2r}.
\end{equation}
With only the Einstein--Hilbert term, vacuum solutions with
constant \(\Lambda_{\mathrm{eff}}\) are Schwarzschild--de Sitter,
\begin{equation}
  \label{eq:SdS}
  f(r)=1-\frac{2M}{r}-\frac{\Lambda_{\mathrm{eff}}}{3}r^2,
\end{equation}
which is singular at \(r=0\) whenever \(M\neq 0\).

\section{Cubic correction in the spherical sector}
\label{sec:ch13-cubic}

Assumption~\ref{ass:ECG} promotes the Einsteinian cubic density
\(\mathcal{P}\) to a term that is solved non-perturbatively in this
sector. Conditional on the first-integral equation of
Section~\ref{sec:status-J} being the static spherical vacuum equation
of the action~\eqref{eq:S-ECG}, that equation is a nonlinear
second-order ODE for \(f\), not an algebraic modification of the mass
function. The mass function below parametrises \(f\); it is not
thereby a solution of that equation. Write
\begin{equation}
  \label{eq:mass-function}
  f(r)=1-\frac{2m(r)}{r}-\frac{\Lambda_{\mathrm{eff}}}{3}r^2.
\end{equation}
The regular-core profile used here is the model
\begin{equation}
  \label{eq:m-reg}
  m(r)=\frac{M r^3}{r^3+r_0^3},
\end{equation}
and therefore
\begin{equation}
  \label{eq:f-reg}
  f(r)
  =1-\frac{2M r^2}{r^3+r_0^3}
   -\frac{\Lambda_{\mathrm{eff}}}{3}r^2.
\end{equation}
The core scale is
\begin{equation}
  \label{eq:r0}
  r_0^3=2\beta_{\mathrm{ECG}}M
\end{equation}
as a phenomenological parametrization (up to the conventional
normalization of \(\mathcal P\)); this relation is not derived from
the action or the first-integral equation. The ECG asymptotic
correction instead defines a characteristic tail length
\(r_*\sim|\alpha M|^{1/5}\), where
\(\alpha=-16\pi G_{\mathrm{eff}}\beta_{\mathrm{ECG}}\) in the stated
action normalization. This does not determine \(r_0\) or identify it
with a core length. The book does not claim a measured value of
\(r_0\).

\paragraph{Limits.}
As \(r\to\infty\), \eqref{eq:f-reg} returns to~\eqref{eq:SdS}.
As \(r\to 0\),
\begin{equation}
  \label{eq:f-core}
  f(r)
  \to
  1-\frac{\Lambda_{\mathrm{eff}}}{3}r^2
  -\frac{2M}{r_0^3}r^2
  +O(r^5),
\end{equation}
which is finite. The curvature scalars remain finite at the origin
for \(r_0\neq 0\). Geodesic completeness through \(r=0\) is a
separate question: the chart~\eqref{eq:sss} is not itself complete,
and a regular origin still requires a smooth matching onto a
neighbourhood covered by local inertial frames of the Levi-Civita
geometry of the composite metric. The frame retains local Lorentz
freedom; the independent Lorentz connection \(B\) is not used here.

\begin{remark}[Model, not an established solution]
  Equation~\eqref{eq:f-reg} is the working profile of the notes.
  For the first-integral equation of Section~\ref{sec:status-J}, it
  is not an exact vacuum solution when \(\Lambda_{\mathrm{eff}}=0\)
  and \(Mr_0\neq0\). The equation also excludes a nonzero integration
  constant for a smooth regular centre for any \(\Lambda_{\mathrm{eff}}\).
  These conclusions are conditional on the reduction and the first
  integral for the invariant \(\mathcal P\) of~\eqref{eq:S-ECG}.
  Whenever a later section
  uses~\eqref{eq:f-reg}, it is using this model, not a theorem that
  ECG produces \(r^3/(r^3+r_0^3)\).
\end{remark}

\section{Horizons and the core}
\label{sec:ch13-horizons}

Roots of \(f(r)=0\) are Killing horizons of~\eqref{eq:sss}. For
\(\Lambda_{\mathrm{eff}}=0\) and \(r_0>0\) there is a range of
\(M\) in which two positive roots exist, one, or none, according
to the ratio \(M/r_0\). For astrophysical masses and microscopic
\(r_0\) the exterior is indistinguishable from Schwarzschild at
the photon sphere and beyond. The core is hidden behind a horizon
if a horizon exists, and is naked if it does not.

NAMCG does not require every compact object to be horizonless.
It requires that the Einstein singularity not be accepted as
fundamental. Whether LIGO ringdowns or EHT images can see
\(r_0\) is a question for Chapter~16, not a claim of this
chapter.

The Einstein-plus-extra-tensor form of the infrared equation
is collected in Appendix~\ref{app:field-eq}.

\section{Weak-field exterior and the extra Poisson equation}
\label{sec:ch13-poisson}

Far from a core, \(r\gg r_0\), the cubic term is negligible and
the geometry is Einstein plus \(\Tent_{\mu\nu}\). In the Newtonian
gauge
\begin{equation}
  \label{eq:newton-gauge}
  ds^2
  =-(1+2\Phi)\,dt^2
   +(1-2\Phi)\delta_{ij}dx^i dx^j,
\end{equation}
the reversible, static piece of the entropic stress reduces to an
effective density
\begin{equation}
  \label{eq:rho-ent}
  \rho_{\mathrm{ent}}
  =T_{00}^{\mathrm{ent,rev}}.
\end{equation}
Poisson's equation becomes
\begin{equation}
  \label{eq:poisson}
  \nabla^2\Phi
  =4\pi\Geff
   \bigl(\rho_{\mathrm{b}}+\rho_{\mathrm{ent}}\bigr),
\end{equation}
where \(\rho_{\mathrm{b}}\) is baryonic density. Conservation,
Theorem~\ref{thm:conservation}, requires that \(\rho_{\mathrm{ent}}\)
not be an arbitrary halo function. It is determined by the
variation of \(U_{\mathrm{ent}}\) for the condensate configuration
sourced by \(\rho_{\mathrm{b}}\).

A useful schematic, sufficient for circular orbits and not
sufficient for lensing tomography, is a source-tracking extra
density
\begin{equation}
  \label{eq:rho-track}
  \rho_{\mathrm{ent}}(r)
  =\frac{1}{4\pi\Geff}
   \nabla\cdot
   \bigl(
     \mu_{\mathrm{ent}}(g_{\mathrm{N}})\,\nabla\Phi_{\mathrm{b}}
   \bigr),
\end{equation}
where \(g_{\mathrm{N}}=|\nabla\Phi_{\mathrm{b}}|\) is the Newtonian
baryonic field and \(\mu_{\mathrm{ent}}\) is a dimensionless
interpolating function fixed by the mutual-information functional.
Equation~\eqref{eq:rho-track} is written to make the structure
visible: \(\rho_{\mathrm{ent}}\) is slaved to baryons. It is not a
six-parameter halo fitted independently of \(\rho_{\mathrm{b}}\).

\begin{assumption}[Reversible extra density tracks baryons]
  \label{ass:track}
  In static, weak-field systems the leading reversible entropic
  density is a local or mildly nonlocal functional of the baryonic
  Newtonian field, of the form~\eqref{eq:rho-track} or a convolution
  thereof with kernel width of order the core scale of the
  constituent condensates.
\end{assumption}

Assumption~\ref{ass:track} is the galactic input. It is narrower
than ``there exists dark matter,'' and it is stronger than anything
proved from \(I(A:B)\) in Chapter~6. Chapter~17 lists a derivation
of \(\mu_{\mathrm{ent}}\) from \(I(A:B)\) as open.

\begin{assumption}[External field enters the interpolator]
  \label{ass:efe}
  For a subsystem whose internal Newtonian field
  \(\mathbf{g}_{\mathrm{N}}\) is embedded in a baryonic field
  \(\mathbf{g}_{\mathrm{ext}}\) that varies slowly across the
  subsystem, the reversible boost is evaluated on the total
  field,
  \begin{equation}
    \label{eq:nu-efe}
    \mathbf{g}
    =\nu_{\sqrt{}}(g_{\mathrm{tot}})\,\mathbf{g}_{\mathrm{N}}
     +\Bigl[\nu_{\sqrt{}}(g_{\mathrm{ext}})
            -\nu_{\sqrt{}}(g_{\mathrm{tot}})\Bigr]
      \mathbf{g}_{\mathrm{ext}},
  \end{equation}
  with \(g_{\mathrm{tot}}=|\mathbf{g}_{\mathrm{N}}+\mathbf{g}_{\mathrm{ext}}|\)
  and \(\nu_{\sqrt{}}\) as in~\eqref{eq:nu-sqrt}.
  In the external-field-dominated limit
  \(g_{\mathrm{N}}\ll g_{\mathrm{ext}}\),
  \begin{equation}
    \label{eq:efe-sat}
    \gamma
    \equiv\frac{g_{\mathrm{int}}}{g_{\mathrm{N}}}
    \to\nu_{\sqrt{}}(g_{\mathrm{ext}}),
  \end{equation}
  not \(\nu_{\sqrt{}}(g_{\mathrm{N}})\).
\end{assumption}

Assumption~\ref{ass:efe} is a coefficient. It is what
\(U_{\mathrm{ent}}[I(A:B)]\) must do if mutual information
sees the Galactic baryons as well as the binary. It is not
implied by Theorem~\ref{thm:conservation}, and it is not a
solution of Problem~H. The Galactic scale at the Sun,
\(g_{\mathrm{gal}}\approx 2.2\times 10^{-10}\,\mathrm{m\,s^{-2}}\),
already satisfies \(g_{\mathrm{gal}}>g_\dagger\), so
\(\nu_{\sqrt{}}(g_{\mathrm{gal}})\approx 1.1\). Isolated-RAR
boosts of order \(4\) are then illegal for solar-neighborhood
wide binaries. A null Gaia result constrains the isolated
reading of Appendix~\ref{app:empirical}, not \(m_T=0\).


\paragraph{Working interpolator.}
Until Problem~H of Chapter~17 produces
\(\mu_{\mathrm{ent}}\) from \(I(A:B)\), galactic comparisons use
\begin{equation}
  \label{eq:nu-sqrt}
  \nu_{\sqrt{}}(g_{\mathrm{N}})
  =\frac{1}{1-\exp\bigl(-\sqrt{g_{\mathrm{N}}/g_s}\bigr)},
\end{equation}
so that \(v_c=V_{\mathrm{b}}\sqrt{\nu_{\sqrt{}}}\). The default
scale is the published SPARC value
\(g_s=g_\dagger=1.20\times 10^{-10}\,\mathrm{m\,s^{-2}}\)
\cite{McGaugh2016}. This \(\nu\) is the RAR interpolator written
as a boost. It is not
\(1/\bigl(1-\alpha\exp(-g_{\mathrm{N}}/g_s)\bigr)\) and it is not
a per-galaxy radial profile. Appendix~\ref{app:empirical}
records the SPARC ranking and retires the older exponential.

\section{Circular velocities}
\label{sec:ch13-vc}

For circular orbits in the midplane,
\begin{equation}
  \label{eq:vc}
  v_c^2(r)
  =r\partial_r\Phi
  =\Geff
   \frac{M_{\mathrm{b}}(<r)+M_{\mathrm{ent}}(<r)}{r},
\end{equation}
with \(M_{\mathrm{ent}}\) the integral of~\eqref{eq:rho-ent}. Two
regimes follow once \(\mu_{\mathrm{ent}}\) interpolates from \(0\)
at large \(g_{\mathrm{N}}\) to a finite value at small \(g_{\mathrm{N}}\).

\paragraph{Interior / high acceleration.}
\(\rho_{\mathrm{ent}}\) is negligible and \(v_c^2\simeq\Geff
M_{\mathrm{b}}/r\). Solar-System and inner-disk data constrain how
small \(\mu_{\mathrm{ent}}\) must be at those accelerations.

\paragraph{Outer / low acceleration.}
If \(M_{\mathrm{ent}}(<r)\) grows approximately linearly with
\(r\), \(v_c\) flattens. That growth is not mysterious once
\eqref{eq:rho-track} is granted: it is the integral of an extra
density that falls more slowly than \(\rho_{\mathrm{b}}\).

On 175 SPARC rotmod files with fixed mass-to-light ratios,
\eqref{eq:nu-sqrt} at frozen \(g_\dagger\) is statistically
equivalent to the published RAR and about \(0.2\%\) behind
simple MOND in \(\chi^2\), while baryons-only is an order of
magnitude worse (Appendix~\ref{app:empirical}). That ranking
is a coefficient check of Assumption~\ref{ass:track}. It is
not a measurement of \(\Geff\), not a derivation from
\(I(A:B)\), and not evidence reserved for the cover.
Whether a single \(\nu_{\sqrt{}}\) survives clusters remains
a Chapter~16 question.

\section{What this is not}
\label{sec:ch13-not}

It is not a proof that dark matter particles are absent. Cluster
lensing, the CMB, and colliding clusters probe \(\Tent_{\mu\nu}\)
in regimes that~\eqref{eq:rho-track} may fail. If those data
require a stress that does not track baryons, NAMCG must either
supply an additional condensate configuration that does, or yield
the dark-matter interpretation in those regimes.

It is not MOND, though the interpolating function
\(\mu_{\mathrm{ent}}\) invites comparison. MOND is a modification
of inertia or of the gravitational field equation at a fixed
acceleration scale \(a_0\). NAMCG obtains an extra conserved
tensor from entanglement. If \(\mu_{\mathrm{ent}}\) happens to
mimic MOND in disks and fails in clusters, that is a phenomenological
coincidence to be measured, not a founding principle.

It is not a replacement for Chapters~3--12. A successful rotation-
curve fit that used an illegal drag, or a non-conserved extra
density, would not support the framework.

\section{Matching core to exterior}
\label{sec:ch13-match}

A compact object of mass \(M\) with core~\eqref{eq:f-reg} matched
to a weak-field exterior~\eqref{eq:poisson} raises no new
consistency issue beyond ordinary matching of \(f\) and \(f'\) at
an intermediate radius. The cubic term is already negligible
there if \(r_0\) is microscopic. The entropic density \(\rho_{\mathrm{ent}}\) of a single compact
object cannot be quantified from the mutual-information tail alone.
Deriving that density and comparing its field with the Newtonian
field requires varying \(U_{\mathrm{ent}}\) for a specified source
configuration. Galactic flattening, if it occurs, is a collective
effect of many baryonic parcels; it is not established here from
the tail of a single compact core.

\section{Exercises}
\label{sec:ch13-exercises}

\begin{enumerate}
  \item \textbf{Einstein limits.}
    Recover~\eqref{eq:SdS} from~\eqref{eq:f-reg} as \(r_0\to 0\)
    at fixed \(r>0\), and as \(r\to\infty\) at fixed \(r_0\).
  \item \textbf{Core curvature.}
    For \(\Lambda_{\mathrm{eff}}=0\), compute \(R\) at \(r=0\)
    from~\eqref{eq:f-core}. Show it is finite for \(r_0\neq 0\).
  \item \textbf{Horizon condition.}
    With \(\Lambda_{\mathrm{eff}}=0\), discuss the number of
    positive roots of \(f(r)=0\) as a function of \(M/r_0\).
    A graphical argument is enough.
  \item \textbf{Poisson split.}
    Starting from~\eqref{eq:poisson} and~\eqref{eq:rho-track},
    write an effective equation for \(\Phi\) in terms of
    \(\Phi_{\mathrm{b}}\) only. What operator acts on
    \(\Phi_{\mathrm{b}}\)?
  \item \textbf{Flat-curve kinematics.}
    Suppose \(M_{\mathrm{b}}+M_{\mathrm{ent}}=\kappa r\) at large
    \(r\). Show \(v_c\) is constant. What density profile
    \(\rho_{\mathrm{ent}}\) produces that mass growth, and is it
    integrable at infinity?
  \item \textbf{Solar-System bound.}
    Estimate an upper bound on \(\mu_{\mathrm{ent}}\) at
    \(g_{\mathrm{N}}\sim 1\,\mathrm{cm\,s^{-2}}\) from the
    requirement that perihelion precession not be spoiled at
    current precision. Order of magnitude is enough.
  \item \textbf{Illegal halo.}
    Write a spherical \(\rho_{\mathrm{halo}}\propto 1/(r^2+r_s^2)\)
    independent of \(\rho_{\mathrm{b}}\). Which assumption of
    this chapter does it violate, and which theorem of
    Chapter~11 does it risk violating if it cannot be obtained
    from \(U_{\mathrm{ent}}\)?
\end{enumerate}

\section{Bridge to Chapter~14}
\label{sec:ch13-bridge}

The extra stress used in this chapter was the reversible,
static piece \(T_{\mu\nu}^{\mathrm{ent,rev}}\). The working
\(d^{-4}\) drag introduced in Chapter~14 is a phenomenological
addition; it is not derived here from time-dependent mutual
information. Its microscopic dynamical origin and compatibility
with a covariantly conserved stress remain unresolved and are
part of Problem~\ref{sec:status-I}. Chapter~15 integrates a
reduced two-body laboratory model. Neither that laboratory nor
the assumed drag substitutes for quadrupole radiation.